% !PRINTED 244
\weberpage{244}
Bedeutet

\weq{3}{P = t^n + p_1 t^{n-1} + p_2 t^{n-2} + \cdots + p_{n-1} t + p_n}

eine nach dem Modul $p$ irreducible Function $n^{\mathrm{ten}}$ Grades, in der wir der Einfachheit halber den Co\"efficienten der höchsten Potenz $t^n$ gleich $1$ annehmen, so lässt sich aus jeder anderen ganzen Function $F(t)$ mit ganzzahligen Co\"efficienten ein Rest $\Phi(t)$ ableiten, dessen Grad niedriger als $n$ ist, indem man (nach \S.\,3 des ersten Bandes)

\weq{4}{F(t) = QP + \Phi(t)}

setzt. $\Phi(t)$ hat hierin ganzzahlige Co\"efficienten, und wir bezeichnen ihre Beziehung zur Function $F(t)$ als eine Congruenz nach dem Modul $P$.

\gesperrt{Es heissen also zwei Functionen $F(t)$, $F_1(t)$, die denselben Rest $\Phi(t)$ haben, congruent nach dem Modul $P$}, was durch die Formel

\[
F(t) \equiv F_1(t)\ (\mathrm{mod}\ P)
\]

ausgedrückt wird. Damit ist gleichbedeutend, dass $F(t) - F_1(t)$ durch $P$ theilbar ist.

Wenn zwei Functionen $F(t)$ und $F_1(t)$ nicht gleichen Rest geben, sondern zwei Reste, die nach dem Modul $p$ congruent sind, so heissen die Functionen $F$, $F_1$ congruent nach dem Doppelmodul $P$, $p$, und man drückt dies durch eine Formel so aus:

\weq{5}{F(t) \equiv F_1(t)\ (\mathrm{mod}\ P, p).}

Für diese Art der Congruenz gilt der Satz:

\gesperrt{1. Das Product zweier Functionen $F(t)$, $F_1(t)$ kann nicht mit Null congruent sein, wenn nicht der eine Factor mit Null congruent ist.}

Der Beweis ergiebt sich aus dem Algorithmus des grössten gemeinschaftlichen Theilers. Ist nämlich $P_1$ eine ganze Function von niedrigerem Grade als $P$, die nicht nach dem Modul $p$ mit Null congruent ist, so kann man eine Reihe von eben solchen Functionen $P_2$, $P_3$, $\ldots$ von abnehmendem Grade, und die Quotienten $Q_1$, $Q_2$, $\ldots$ gleichfalls als ganze Functionen so bestimmen, dass

\weq{6}{P \equiv Q_1 P_1 + P_2,\quad P_1 \equiv Q_2 P_2 + P_3, \ldots,\quad P_{\nu-2} \equiv Q_{\nu-1} P_{\nu-1} + P_\nu\ (\mathrm{mod}\ p)}

wird, und die Reihe dieser Gleichungen bricht ab, wenn $P_\nu$ eine Constante (ganze Zahl) geworden ist. Diese Constante $P_\nu$ kann
%CONT>>
